\subsection{\resultrevision{Comparison Design}}
\begingroup
\color{resultrevision}
The main comparison includes MPSLGP, PSLGP and NSGPII. PSLGP denotes the IGD-trained PSLGP-I variant~\cite{xu2026pslgp}. All methods are evaluated using generation-49 outputs on the same held-out scenarios, with test seed 968356 and 30 simulation replications. MPSLGP and PSLGP each supply one fixed conditioned heuristic evaluated across 200 preferences; NSGPII supplies its complete final population Pareto front. The comparison comprises 540 MPSLGP task/run tests and 300 tests for each independent method.

MPSLGP allocates 1,000 individuals across each two- or three-task group; each independent baseline task has 1,000 individuals. MPSLGP and NSGPII run for 50 evaluated generations, while PSLGP's 51-generation training supplies its generation-49 best individual. MPSLGP selects the first training-ranked individual for each task. The comparison assesses the quality and cost of these joint and independent training configurations, retaining their native operators and output representations.

The no-transfer ablation uses the same groups, seeds, task quotas, crossover, IGD aggregation, surrogate procedure and test configuration as Full. Its adaptive transfer is disabled and its cross-task probability and budget are zero. Full and no-transfer are evaluated under common reference fronts pooled over these two profiles.

Quality measures are reported as means and standard deviations calculated over all 30 training runs. For each target and indicator, an exact two-sided Wilcoxon signed-rank test pairs outcomes by training run ID. Zero differences are omitted, absolute ties receive midranks, and raw $p\leq0.05$ defines significance without multiplicity adjustment. The statistical unit is the training seed. Arrows follow the favourable signed-rank direction, while bold entries identify the best means. Inference is per comparison; the no-transfer ablation assesses the combined transfer mechanism rather than its individual components.
\endgroup